\documentclass[10pt,a4paper]{amsart}
\usepackage[margin=19mm]{geometry}
\usepackage{amsmath,amssymb,mathtools}
\usepackage[scr=rsfs,cal=euler]{mathalfa}
\usepackage{xfrac}
\usepackage[unicode,hidelinks,bookmarks=false]{hyperref}
\newcommand{\ie}{i.e.\ }
\newcommand{\cf}{cf.\ }
\newcommand{\resp}{respectively\ }
\newcommand{\Subsection}{\subsection}
\providecommand{\nobreakdash}{\nobreak\hbox{-}}
\setlength{\emergencystretch}{2em}
\multlinegap0pt
\usepackage{amssymb}
\newtheorem{lemma}{Lemma}
\newtheorem{theorem}{Theorem}
\title{Weighted $H^p\!- \! L^q$ boundedness of integral operators with rough kernels}
\author{Andrea L. Gallo, M. Silvina Riveros, Lucas A. Vallejos}
\date{Source: arXiv:2607.03567v1 (CC BY 4.0)}
\begin{document}
\maketitle

\noindent\emph{Source and licence.} Weighted $H^p\!- \! L^q$ boundedness of integral operators with rough kernels. Version arXiv:2607.03567v1 (CC BY 4.0), distributed by arXiv under the Creative Commons Attribution 4.0 International licence, http://creativecommons.org/licenses/by/4.0/, as recorded in the arXiv licence metadata for this exact version and shown on the abstract page https://arxiv.org/abs/2607.03567v1. Full licence notice: \texttt{licenses/arxiv-2607.03567-CC-BY-4.0.txt}.\par\smallskip

\noindent\textit{Editorial scope.} Counted unit: the article's theorem Extension (source0.tex lines 452--454). The article's complete original proof of it is delivered with it (source0.tex lines 456--491, one \texttt{proof} environment of 36 lines). No mathematics is written, repaired or re-derived: the statement and the proof are the article's own lines, in the article's own notation, and no retained line is substituted or re-flowed.  Retained uncounted as the source-local context the proof needs: lines 438--439; lines 440--442.  Nothing was omitted from any retained span, so the package carries no omission record.  The retained lines cite 1 works, whose entries are transcribed verbatim from the article's own bibliography source in the e-print and delivered with short editorial labels recorded in the item's reference mapping.  No retained line contains a figure environment, an image-inclusion command or drawing code, so no image is delivered and the package declares no asset dependency.  Editorial formatting only, in addition: an AMS \texttt{amsart} preamble that restates the article's own declarations and macros used by the retained lines, namely the environments \texttt{lemma}, \texttt{theorem} on separate counters, together with the standard packages those lines need.  The retained environments print in this document as Lemma 1, Theorem 1 in order of appearance, whereas the article prints the same items with the numbering of its own full document.  The complete native source of the article as distributed in the arXiv e-print for this exact version is bundled unchanged as \texttt{sources/204400/source0.tex}.
	\begin{lemma}[\cite{Rocha}]\label{Lemma8} Let $w \in \mathcal{A}_{\infty}$ with critical index $\tilde{q}_w$ and critical index $r_w$ for the reverse H\"{o}lder's condition. If $a(\cdot)$ is a $w$-$(p,p_0,d)$ atom, then $a(\cdot) \in H^{p}_w(\mathbb{R}^{n})$. Moreover, there exists a positive constant $C$ independent of the atom $a$ such that $||a||_{H^{p}_w} \leq C$. 
	\end{lemma}

	\begin{theorem}[\cite{Rocha}]\label{Theorem9}
		Let $f \in \hat{\mathcal{D}_0}$, and $0<p \leq 1$. If $w \in \mathcal{A}_{\infty}$, then there exist a sequence of $w$-$(p,p_0,d)$ atoms $\{a_j\}$ and a sequence of scalar $\{\lambda_j\}$ with $\sum_j |\lambda_j|^{p} \leq c||f||^{p}_{H^{p}_w}$ such that $f=\sum_j \lambda_ja_j$, where the convergence is both in $L^{s}(\mathbb{R}^{n})$ and pointwise, for each $1<s<\infty$.
	\end{theorem}

	\begin{theorem}\label{Extension}
		Let $T$ be a bounded linear operator from $L^{p_0}(\mathbb{R}^{n})$ to $L^{q_0}(\mathbb{R}^{n})$ for some $1<p_0<\infty$ and $p_0<q_0<\infty$. Let $w \in \mathcal{A}_{\infty}$ with critical index $r_w$, $0<p \leq \min \{1, \frac{r_w-1}{r_w} p_0 \}$, and let $q\geq p$ be a real number. Then $T$ can be extended to a bounded linear operator $H^{p}_w(\mathbb{R}^{n})-L^{q}_w(\mathbb{R}^{n})$ if and only if $Ta$ is uniformly bounded into the $L^{q}_w$ norm for all $w-(p,p_0,d)$ atoms $a$.
	\end{theorem}

	\begin{proof}
		By Lemma \ref{Lemma8}, we have that $a \in H^{p}_w(\mathbb{R}^{n})$ for all $w$-$(p,p_0,d)$ atom. Then, if $T$ can be extended to an $H^{p}_w(\mathbb{R}^{n})-L^{q}_w(\mathbb{R}^{n})$, we obtain that $||Ta||_{L^{q}_w} \leq c_p||a||_{H^{p}_w}$. Therefore, by Lemma \ref{Lemma8} we have that $||Ta||_{L^{q}_w} \leq c_p$ for all $a$ $w$-$(p.p_0,d)$ atom.
		
		Conversely, by Theorem \ref{Theorem9}, given $f \in \hat{\mathcal{D}_0}$ there exist a sequence of $w$-$(p,p_0,d)$ atoms $\{a_j\}$ and a sequence of scalar $\{\lambda_j\}$ with $\sum_j |\lambda_j|^{p} \leq c||f||^{p}_{H^{p}_w}$ such that $f=\sum_j \lambda_ja_j$, where the convergence is both in $L^{p_0}(\mathbb{R}^{n})$ and pointwise. 
		Since $T$ is a bounded operator from $L^{p_0}(\mathbb{R}^{n})$ in $L^{q_0}(\mathbb{R}^{n})$, we have that $\sum_j \lambda_jTa_j$ converges to $Tf$ in $L^{q_0}(\mathbb{R}^{n})$.
		Then, there exists a subsequence $\{r_n\}$ such that $\lim_{n \to \infty} \sum_{j=1}^{r_n} \lambda_j Ta_j(x)=Tf(x)$ a.e.$x \in \mathbb{R}^{n}$.
		Then, 
		\begin{equation}\label{Tpuntual}
			|Tf(x)| \leq \sum_j |\lambda_jTa_j(x)|, \quad a.e.x \in \mathbb{R}^{n}.
		\end{equation}
		
		In the follow we analyze the cases $p\leq q<1$ and $q\geq 1$. First, if $p\leq q<1$, \eqref{Tpuntual} implies that 
		$$ \int_{\mathbb{R}^{n}} |Tf(x)|^{q} w(x) dx \leq \sum_j |\lambda_j|^{q} \int_{\mathbb{R}^{n}} |Ta_j(x)|^{q} w(x) dx.$$
		If we suppose that $||Ta||_{L^{q}_w} \leq C_p$ for all $w$-$(p.p_0,d)$ atom $a$, we have that 
		$$ \Big(\int_{\mathbb{R}^{n}} |Tf(x)|^{q} w(x) dx \Big)^{\frac{1}{q}} \leq C_p \Big(\sum_j |\lambda_j|^{q}\Big)^{\frac{1}{q}}.$$
		It is easy see that $ \left( \sum_j |\lambda_j|^q \right)^{\frac{1}{q}} \leq \Big( \sum_j|\lambda_j|^{p}\Big)^{\frac{1}{p}} \leq c ||f||_{H^{p}_w}$, 
		so we have that 
		$$ \Big(\int_{\mathbb{R}^{n}} |Tf(x)|^{q} w(x) dx \Big)^{\frac{1}{q}} \leq C_p \Big(\sum_j |\lambda_j|^{q}\Big)^{\frac{1}{q}}  \leq C_p ||f||_{H^{p}_w},$$
		
		for all $f \in \hat{\mathcal{D}_0}$, and the Theorem follows by the density of $\hat{\mathcal{D}_0}$ in $H^{p}_w(\mathbb{R}^{n})$.
		
		Now, if $q\geq 1$, by Minkowsky's inequality, monotone convergence theorem and \eqref{Tpuntual} implies that 
		$$
		\begin{aligned}
			\left(\int_{\mathbb{R}^{n}} |Tf(x)|^{q} w(x) dx\right)^{\frac{1}{q}} &\leq \left(\int_{\mathbb{R}^{n}} \left(\sum_{j=1}^{\infty} |\lambda_j| |Ta_j(x)|\right)^{q} w(x) dx\right)^{\frac{1}{q}}\\
			&\leq  \lim_{n \to \infty} \left(\int_{\mathbb{R}^{n}} \left(\sum_{j=1}^{n} |\lambda_j| |Ta_j(x)|\right)^{q} w(x) dx\right)^{\frac{1}{q}}\\
			&\leq  \lim_{n \to \infty} \Big\| \sum_{j=1}^{n} |\lambda_j| |Ta_j| \Big\|_{L^{q}(w)}\\
			&\leq \lim_{n \to \infty} \sum_{j=1}^{n}|\lambda_j| \ || Ta_j||_{L^{q}(w)}\\
			&\leq C \lim_{n \to \infty} \sum_{j=1}^{n}|\lambda_j|.
		\end{aligned}
		$$
		
		Since $0<p\leq 1$ , we have that $\sum_j |\lambda_j| \leq \Big( \sum_j|\lambda_j|^{p}\Big)^{\frac{1}{p}} \leq c ||f||_{H^{p}_w}$, and then  we have that 
		$$ \Big(\int_{\mathbb{R}^{n}} |Tf(x)|^{q} w(x) dx \Big)^{\frac{1}{q}} \leq C_p ||f||_{H^{p}_w},$$
		for all $f \in \hat{\mathcal{D}_0}$, and the Theorem follows by the density of $\hat{\mathcal{D}_0}$ in $H^{p}_w(\mathbb{R}^{n})$.
	\end{proof}

\begin{thebibliography}{99}
\bibitem[Rocha]{Rocha} Rocha P., On the atomic and molecular decomposition of weighted Hardy spaces. Rev. Uni\'on Mat. Argent. \textbf{61(2)}, 229--247, (2020).
\end{thebibliography}
\end{document}
